\subsection{Homepage}
La homepage mantiene la struttura originale: suddivisa in header, sezione di benvenuto, funzionalità, prezzi e footer. L'obiettivo principale è rimasto quello di presentare la piattaforma ai nuovi utenti e invitarli a registrarsi. In questa fase si è scelto di non stravolgere l'impostazione iniziale, dato che la homepage non è emersa come la sezione più critica del sito e non presentava tanti problemi di usabilità. L'intervento si è quindi limitato alla risoluzione dei problemi riscontrati, senza spostare il focus dalle aree del sistema che richiedevano maggiore attenzione.

\paragraph{Design pattern applicati}

\begin{itemize}
    \item \textbf{Sign-In Tools} (Tidwell): Il pattern prevede di posizionare nella parte superiore destra dell'interfaccia i controlli di accesso, come i pulsanti di login e registrazione, in quanto è la posizione convenzionale in cui gli utenti si aspettano di trovarli. Nella versione precedente il pulsante di accesso "Accedi" aveva un aspetto meno evidente rispetto a "Inizia Subito", risultando quasi non cliccabile, e i due pulsanti non comunicavano chiaramente la distinzione tra accesso e registrazione. Nella versione riprogettata entrambi i pulsanti sono stati uniformati graficamente e rinominati in "Accedi" e "Registrati", rendendone immediatamente chiaro il ruolo e la differenza.

    Problemi di usabilità risolti:
    \begin{itemize}
        \item P1 (Il pulsante "Inizia Subito" dovrebbe essere "Registrati" per coerenza con "Accedi"): Il pulsante è stato rinominato \textit{"Registrati"}.
        \item P81 (Il pulsante "Accedi" è meno evidente e può sembrare non cliccabile): I due pulsanti sono stati uniformati graficamente e accorpati in alto a destra.
    \end{itemize}

\end{itemize}

\paragraph{Altre modifiche}
Sono state inoltre apportate diverse correzioni mirate a migliorare la coerenza e l'usabilità dell'interfaccia nella homepage. In particolare, è stata eliminata la duplicazione del pulsante \textit{Inizia Subito}, mantenendo un'unica call-to-action principale per evitare ambiguità. La sezione dei piani di abbonamento è stata resa più visibile e uniforme, seguendo l'idea dei mockup, migliorando il contrasto e rimuovendo lo sfondo trasparente per una maggiore leggibilità. Tutti i pulsanti associati ai piani sono stati resi coerenti sia nel testo sia nel comportamento: le etichette sono state uniformate in \textit{Registrati} e ogni bottone risulta ora correttamente cliccabile. Infine, nel footer sono stati rimossi i collegamenti non funzionanti sotto le sezioni \textit{Azienda} e \textit{Legale}, lasciando solo i link attivi a \textit{Funzionalità} e \textit{Prezzi}, che indirizzano correttamente alle rispettive sezioni della pagina. \\

Problemi di usabilità risolti:
\begin{itemize}
    \item P2 (Duplicazione del pulsante "Inizia Subito"): È stato rimosso il pulsante centrale, mantenendo un'unica call-to-action in alto a destra.
    \item P4 (Le sezioni dei piani di abbonamento non sono ben visibili): Le card sono state uniformate e rese più evidenti, eliminando lo sfondo trasparente.
    \item P5 (Il pulsante "Get Started" nei piani non è cliccabile): Tutti i pulsanti sono stati resi cliccabili e rinominati \textit{Inizia ora}.
    \item P6 (Il piano "Standard" ha un pulsante non cliccabile): Il pulsante è ora attivo e coerente con gli altri (\textit{Inizia ora}).
    \item P7 (Testo incoerente del pulsante "Inizia Subito, Gratis"): Il testo è stato uniformato a \textit{Registrati}.
    \item P8 (Link non funzionanti nel footer): Rimossi i link non attivi e mantenuti solo quelli funzionanti.
\end{itemize}